% !TeX root = induk.tex
\chapter{\normalsize KESIMPULAN DAN SARAN}
\raggedbottom

Bab ini menyimpulkan hasil temuan utama dari analisis STATIS dan Dynamic Biplot mengenai dinamika pembangunan ekonomi regional di Provinsi Banten, DKI Jakarta, dan Jawa Barat, serta memberikan saran rekomendasi kebijakan maupun pengembangan penelitian selanjutnya.

%-----------------------------------------------------------------------------%
% --- 5.1 ---
\section{\normalsize Kesimpulan}
%-----------------------------------------------------------------------------%
Berdasarkan hasil analisis dan pembahasan yang telah dilakukan, dapat ditarik beberapa kesimpulan sebagai berikut:

\begin{enumerate}
    \item \textbf{Karakteristik Struktur Kompromi (Menjawab Tujuan 1):}
    Melalui sintesis matriks kompromi STATIS (Skenario 1 $\mathbf{W}_c$ dan Skenario 2 $\mathbf{S}_c$), diperoleh representasi rata-rata spasio-temporal yang stabil selama periode 2020--2024. Dua dimensi utama berhasil diidentifikasi: Dimensi~1 merepresentasikan ``Taraf Ekonomi dan Kualitas Sumber Daya Manusia'' (didominasi oleh kontribusi positif dari indikator PDRB, IPM, RLS, dan IPG) serta menjelaskan bagian terbesar variansi total (\AppDimOnePct{}\%). Dimensi~2 merepresentasikan ``Kesehatan dan Ketahanan Sosial'' (didominasi oleh Angka Harapan Hidup L/P dan PKKP). Vektor bobot interstruktur ($\boldsymbol{\alpha}_W$ dan $\boldsymbol{\alpha}_S$) terdistribusi secara seimbang di antara kelima tahun pengamatan, yang membuktikan representativitas yang adil tanpa distorsi tahun ekstrem.

    \item \textbf{Signifikansi Stabilitas Temporal dan Vektor Struktur (Menjawab Tujuan 2):}
    Analisis interstruktur STATIS menghasilkan koefisien RV yang sangat tinggi (seluruhnya di atas $\AppRVMin{}$), membuktikan bahwa hierarki kemiripan antar-wilayah bersifat persisten lintas waktu. Uji permutasi Monte Carlo ($B = 1000$) membuktikan bahwa struktur kompromi tersebut \textbf{signifikan secara statistik} ($p\text{-value} = \AppPValuePerm{} < 0{,}05$) dan bukan merupakan hasil fluktuasi acak. Kompresi interstruktur ke dalam vektor struktur $\tilde{\boldsymbol{\beta}} = \lambda_1 \boldsymbol{\alpha}_1$ secara longitudinal merekam dampak krisis pandemi (2020--2021) dan pemulihan (2022--2024) yang nyata, sementara analisis transversal membuktikan adanya heterogenitas respons temporal yang unik antar-provinsi.

    \item \textbf{Pola Lintasan Trajektori Spasio-Temporal (Menjawab Tujuan 3):}
    Visualisasi \textit{Dynamic Biplot} berhasil menggambarkan lintasan (\textit{trajectory}) pergerakan 41 kabupaten/kota secara dinamis dari masa krisis (2020) hingga masa pemulihan (2024). Tiga pola pergerakan diidentifikasi: (1) pola pemulihan berbentuk huruf V (\textit{V-shape rebound}) pada kota-kota metropolitan berbasis jasa (seperti DKI Jakarta, Kota Bandung, Kota Tangerang Selatan), (2) pola progresif bertahap (\textit{monotone improvement}) pada wilayah penyangga, dan (3) pola stagnasi (\textit{inertia}) pada kabupaten pedalaman bagian selatan (seperti Lebak, Pandeglang, Garut) yang respons pemulihannya lambat.

    \item \textbf{Pengelompokan Spasial dan Dinamika Transisi (Menjawab Tujuan 4):}
    Penerapan algoritma K-Means ($k=3$) pada skor kompromi berhasil mengklasifikasikan wilayah secara objektif menjadi tiga kelompok: Klaster~1 (Tertinggal/Banten Selatan dan pedalaman), Klaster~2 (Berkembang/Wilayah penyangga), dan Klaster~3 (Maju/Metropolitan). Matriks transisi temporal lintas tahun mengonfirmasi dominasi diagonal utama yang sangat tinggi (stabilitas pengelompokan yang kokoh), namun tetap menangkap adanya mobilitas naik (\textit{upward migration}) yang selektif pada wilayah tertentu (seperti Kabupaten Kepulauan Seribu yang bermigrasi dari Klaster 1 ke Klaster 2 pada tahun 2023).

    \item \textbf{Pengembangan Dashboard Eksplorasi Interaktif (Menjawab Tujuan 5):}
    Aplikasi berbasis web dikembangkan menggunakan \textit{framework} Dash dan Plotly sebagai instrumen eksplorasi visual interaktif yang andal. Fitur-fitur seperti pemetaan spasial dinamis, visualisasi trajektori, seleksi tahun, dan rotasi sumbu \textit{baseline} berhasil memfasilitasi pemantauan indikator pembangunan secara interaktif oleh peneliti dan pemangku kebijakan. Keandalan model terkonfirmasi melalui validasi eksternal (tabulasi silang) yang menunjukkan kesesuaian yang sangat tinggi antara klaster STATIS dengan klasifikasi normatif kategori IPM UNDP.
\end{enumerate}

%-----------------------------------------------------------------------------%
% --- 5.2 ---
\section{\normalsize Saran}
%-----------------------------------------------------------------------------%
Berdasarkan kesimpulan di atas, diajukan beberapa saran sebagai berikut:

\begin{enumerate}
    \item \textbf{Implikasi Kebijakan \textit{Targeted} Regional:}
    Mengingat struktur ketimpangan yang persisten dan terpola secara spasial, kebijakan pemerataan pembangunan tidak cukup hanya mengandalkan pertumbuhan ekonomi nasional. Diperlukan intervensi struktural yang \textit{targeted} pada wilayah Klaster~1 (Tertinggal), khususnya Kabupaten Lebak, Pandeglang, dan Garut, dengan fokus pada peningkatan kapabilitas dasar (RLS, AHH) serta pengurangan kemiskinan melalui strategi \textit{regional spillover} dari pusat pertumbuhan DKI Jakarta.
    
    \item \textbf{Pengembangan Penelitian Lanjutan:}
    Penelitian ini dibatasi pada indikator makroekonomi tahunan dari 3 provinsi. Penelitian selanjutnya disarankan untuk menggunakan data dengan resolusi waktu yang lebih tinggi (triwulanan) apabila tersedia, memperluas cakupan wilayah ke seluruh Pulau Jawa atau Indonesia dengan mengintegrasikan metode \textit{Dual STATIS} untuk analisis korelasi indikator, serta menambahkan indikator sosial-infrastruktur guna memperkaya dimensi analisis.

    \item \textbf{Pengembangan Model Prediktif Spasio-Temporal (\textit{Forecasting}):}
    Karena penelitian ini terbatas pada analisis deskriptif-eksploratif data historis, penelitian selanjutnya sangat disarankan untuk mengintegrasikan model prediktif (peramalan) guna memproyeksikan skor wilayah ($F_k$) dan loading indikator ($G_k$) ke tahun berikutnya ($k = K+1$). Hal ini dapat diwujudkan dengan menerapkan pemodelan deret waktu (seperti ARIMA, Vector Autoregression (VAR), atau model \textit{Grey Systems} GM(1,1) yang sangat handal untuk deret data pendek $K=5$) langsung pada barisan koordinat kompromi ($\mathbf{f}_{i,k}$ dan $\mathbf{g}_{j,k}$) dan lintasan komponen vektor struktur ($\tilde{\beta}_k$). Dengan memproyeksikan lintasan koordinat ini, pemangku kebijakan dapat mengantisipasi pergeseran status kesejahteraan wilayah dan dinamika hubungan antar-indikator secara prospektif.

    \item \textbf{Pengembangan Platform \textit{Dashboard} Interaktif:}
    Dashboard berbasis web (Dash/Plotly) yang dikembangkan dalam penelitian ini dapat dikembangkan lebih lanjut dengan menambahkan fitur pembaruan data secara otomatis dari API BPS Open Data, sehingga pemantauan dinamika pembangunan regional dapat dilakukan secara \textit{near real-time} oleh pemangku kepentingan kebijakan.
\end{enumerate}

